\title{opensw 說明文件}
\maketitle

\section{開始}

opensw 是照 Scientific Word 5.5 重做的數學編輯器。選單、快捷鍵、檔案格式盡量跟 5.5 一樣。

\begin{itemize}
\item 桌面版（Windows）能做全部的事：開檔存檔、排版 PDF、計算、AI 控制。
\item 網頁版可以打字和計算，但不能直接讀寫電腦上的檔案，也不能排版。
\item 按 F1 看快捷鍵表。這份說明可以用 編輯 \TEXTsymbol{>} 尋找（Ctrl+Q）找字。
\end{itemize}

\subsection{畫面}

\includegraphics{manual/zh-TW/window.png}

\emph{開著 檔案 > 範例 > 數學範例 的畫面。}

\begin{itemize}
\item 最上面：選單和工具列。檢視 \TEXTsymbol{>} 工具列 可以開關每一列。
\item 右邊：插入註解、標記、交叉引用、連結、索引項目、TeX 和 HTML 欄位、考卷公式的按鈕，還有往前、往後。
\item 最下面：標籤列（游標所在地方的項目標籤、段落標籤、文字標籤），下面是狀態列（文字或數學、縮放比例、AI 在做什麼）。
\item 公式在畫面上是紅色，工具 \TEXTsymbol{>} 使用者設定 可以改顏色；排版出來是黑色。
\end{itemize}

\section{檔案}

\includegraphics{manual/zh-TW/menu-file.png}

\emph{檔案選單。}

\subsection{開檔和存檔}

\begin{itemize}
\item 開啟：6.x 的 .sci、5.x 的 .tex、考卷範本 .qiz 都可以。
\item 儲存：主要格式是 5.x 的 .tex（5.5 打得開，中文用 Big5）。打開的 .tex 會存回原檔，存之前留一份 .bak；沒改過就不重寫。從 .sci 打開的文件照樣存成 .sci。
\item 另存成 .sci：6.x 的格式，給需要的人。
\item 覆蓋舊檔前會留一份 .bak（工具 \TEXTsymbol{>} 使用者設定 可以關掉）。
\item 開啟網址：網址結尾是 .tex 或 .sci 就在這裡開；其他網址用瀏覽器開。網址後面加「 as .tex」可以把沒有 .tex 結尾的檔當文件開。
\item 最近的檔案（桌面版）、範例、文件資訊（字數、公式數、圖片數）。
\end{itemize}

\subsection{片段}

片段是存起來的一小段文字和公式，隨時可以插進文件。跟 5.5 的 .frg 檔一樣，內容是 5.x 的 LaTeX。

\begin{itemize}
\item 存成片段：先選取內容，再取名字。
\item 匯入片段：從清單挑一個插到游標的地方。也可以從 5.5 的 Frags 資料夾加入 .frg 檔。
\item 工具列的「片段…」清單一樣可以插入。內建幾個物理常數。
\item 跟 5.5 一樣，按住 Ctrl 打片段的名字（兩個以上的英文字母或數字，不含資料夾），放開 Ctrl 就插入。名字不能用 C、V、X、A、Z、Y 開頭：這幾個留給複製、貼上、剪下、全選、復原、重做。
\item 片段存在這台電腦，所有文件共用。
\end{itemize}

\subsection{匯入和匯出}

\begin{itemize}
\item 匯入圖片：PNG、JPG、GIF、SVG 等。
\item 匯出 LaTeX、HTML、Word（.docx）、給 AI 讀的 Markdown。
\item 匯出 Word 時公式會變成 Word 的方程式，有編號的公式編號靠右。
\end{itemize}

\subsection{版面設定、預覽、排版}

\begin{itemize}
\item 版面設定：邊界、頁首、頁尾、頁碼。只用在「預覽」，不影響排版 PDF（跟 5.5 一樣）。
\item 預覽（不排版）：照畫面的樣子分頁，送到系統的列印預覽，可以直接印或存 PDF。
\item 排版 PDF（Ctrl+P，桌面版）：用 LaTeX 排成 PDF。5.x 的文件會用 5.x 的巨集排，看起來跟 5.5 排的一樣。
\end{itemize}

\section{編輯}

\begin{itemize}
\item 復原、重做、剪下、複製、貼上、全選。
\item 選擇性貼上：當純文字、當 LaTeX（節標題、公式照 LaTeX 讀）、或當一條公式。
\item 複製成圖片：游標在公式裡就複製那條公式；選取圖片就複製那張圖。解析度在 工具 \TEXTsymbol{>} 匯出設定。
\item 屬性（Ctrl+F5）：公式改 TeX、表格改對齊和格線、圖片改大小和標題、繪圖改公式和範圍。
\item 表格：插入或刪除列和欄、合併和分割儲存格。在表格裡按 Tab 跳下一格。
\item 尋找（Ctrl+Q）、取代（Ctrl+W）：文字和公式裡都找。
\end{itemize}

\includegraphics{manual/zh-TW/menu-edit.png}

\emph{編輯選單。}

\includegraphics{manual/zh-TW/props.png}

\emph{公式的屬性（游標在公式裡按 Ctrl+F5）：TeX、編號、標記。}

\section{插入}

\begin{itemize}
\item 數學（Ctrl+M 或打 \$）開始一條公式；獨立公式 Ctrl+D。
\item 分數 Ctrl+F、根號 Ctrl+R、上標 Ctrl+H、下標 Ctrl+L。
\item 運算子、括號、矩陣、函數名稱、二項式、裝飾（帽子、橫線、箭頭）。
\item 表格、表格標題。
\item 標記和交叉引用：先在要引用的地方放標記，再插入交叉引用。
\item 註解：註腳、邊註。點小數字可以改內容。
\item 引用文獻、索引項目、超連結。
\item TeX 欄位：直接寫 LaTeX，畫面顯示灰底，排版時照用。
\item HTML 欄位：只在匯出 HTML 時用。
\item 考卷公式：產生考卷時才算的公式（見「考卷產生器」）。
\item 間距、單位名稱。
\end{itemize}

\includegraphics{manual/zh-TW/menu-insert.png}

\emph{插入選單。}

\includegraphics{manual/zh-TW/matrix.png}

\emph{插入 > 矩陣：列數、欄數、括號。}

\section{檢視、前往、標籤}

\begin{itemize}
\item 檢視：放大縮小；顯示或隱藏標記、索引項目、輔助線、TeX 欄位；工具列（每一列都可以開關，有「重設」）。
\item 檢視 \TEXTsymbol{>} 改動與 LaTeX 源碼（Ctrl+Shift+D；5.5 沒有，不打開就不會出現）：左邊是打開時的原檔，只能看；右邊三個分頁。工作版本：之後改過的區塊上色，AI 改的是紫色、自己改的是黃色。LaTeX 源碼：存檔會寫出的內容，只能看、隨時更新；重新寫過的部分有標記，其他是原檔的文字。改動與錯誤：每一條改動可以接受（保留、不再標出）或退回（換回原檔）；上次排版的錯誤翻成白話，並指出在哪一塊；推導檢查：獨立公式裡每一步 $a=b$（下一行開頭的 $=b$ 也算）會判定為確定相等、確定不相等（附上兩邊不同的點），或無法判斷（這時不說錯）。
\item 前往：上一節、下一節、標記、連結（游標在連結上）、上一個位置、下一個位置。
\item 標籤：F2 到 F12 套用標籤（標籤 \TEXTsymbol{>} 功能鍵 可以自己改）；項目標籤（清單、定理類）、段落標籤（節、小節）、文字標籤（粗體、強調、字級）。
\item 標籤 \TEXTsymbol{>} 外觀：改每個標籤在畫面上的字型、大小、顏色，不影響排版。
\end{itemize}

\includegraphics{manual/zh-TW/menu-view.png}

\emph{檢視選單。}

\includegraphics{manual/zh-TW/review.png}

\emph{加了一句話之後的 檢視 > 改動與 LaTeX 源碼：改過的段落兩邊都有底色。}

\includegraphics{manual/zh-TW/review-source.png}

\emph{LaTeX 源碼分頁：存檔會寫出的內容，在這裡只能看不能改。}

\includegraphics{manual/zh-TW/review-list.png}

\emph{改動與錯誤分頁：每條改動可以接受或退回；排版錯誤；檢查推導的每一步。}

\includegraphics{manual/zh-TW/menu-tag.png}

\emph{標籤選單。}

\includegraphics{manual/zh-TW/tag-appearance.png}

\emph{標籤 > 外觀：每個標籤在畫面上的樣子。}

\section{計算}

游標放在公式裡，從計算選單選一個動作。結果接在公式後面。計算用的是 SymPy，第一次要載入幾秒；桌面版在 工具 > 計算引擎 可以改用這台電腦裝的 Mathematica。計算是 opensw 自己的功能：5.5 用的是 MuPAD，opensw 沒有，所以答案的寫法可能和 5.5 不同（值相等），有些指令給的答案也可能多一個或少一個。

\includegraphics{manual/zh-TW/menu-compute.png}

\emph{計算選單，打開了 微積分 子選單。}

\begin{itemize}
\item 計算（Ctrl+E）、數值計算、化簡、因式分解、展開、合併、改寫。
\item 檢查是否相等：看等式兩邊是不是相等。
\item 解：精確、數值、整數、遞迴。解常微分方程：精確、數值。
\item 解常微分方程 > 數值：解一個起始值問題，打成一欄的矩陣，一列一條方程或一個起始值（例如 $y'=-y$ 和 $y(0)=1$；任何階的方程或方程組：$y$ 的 $n$ 階方程要有起始值 $y(x_0)$、$y'(x_0)$、\ldots、$y^{(n-1)}(x_0)$，例如 $y''+y=0$、$y(0)=1$、$y'(0)=0$）。和 5.5 一樣回答「已定義的函數：$y$」：解變成函數，之後計算 $y(1)$ 會得到數字，$y$ 作用在一欄數字上得到一欄函數值，畫 $y$（繪圖 2D > 直角座標或常微分方程）就畫出這個解。
\item 多項式：部分分式、依變數整理。
\item 微積分：微分、不定積分、分部積分、換變數、近似積分、畫近似積分圖、冪級數、Padé 近似、求極值、迭代、隱函數微分。
\item 向量微積分：梯度、散度、旋度、拉普拉斯、黑塞矩陣、雅可比矩陣、朗斯基行列式、位勢、向量位勢。基底變數預設是 $x,y,z$；向量微積分 > 設定基底變數（或 工具 > 計算引擎）可以改，計算裡的 $\nabla$ 也用它：設成 $u,v,w$ 時，$\nabla (cuv)$ 是 $(cv,cu,0)$。
\item 單形法：最大化、最小化、可行性、標準化、對偶。打成一行矩陣，第一列是目標函數，其他列是限制式，變數都當成 $\geq 0$。
\item 統計、拉氏轉換、傅立葉轉換、矩陣（行列式、反矩陣、特徵值、秩、零空間⋯⋯）。
\item 定義：例如 $f(x)=x^{2}$ 選「新定義」，之後的計算都會用到。定義帶字母下標的名稱（例如 $a_{i}=3i$）時，會和 5.5 一樣問下標怎麼解讀：函數的引數（這時 $a_{2}$ 是 $6$，右邊的 $i$ 是引數，不是虛數單位），或名稱的一部分（$a_{2}$ 還是 $a_{2}$）。
\item 用 Lean 檢查：把等式交給 Lean 4 證明（要先在 工具 \TEXTsymbol{>} Lean 設定 選 Lean 專案資料夾）。
\end{itemize}

\includegraphics{manual/zh-TW/compute.png}

\emph{按 Ctrl+E 答案接在公式後面；解方程把解寫在後面。}

\subsection{繪圖}

游標放在公式裡，選 計算 \TEXTsymbol{>} 繪圖 2D 或 繪圖 3D。圖畫好後選取它按 Ctrl+F5 可以改公式、範圍、格線數。5.5 檔裡的繪圖會照原本的設定重新畫。

\begin{itemize}
\item 2D 直角座標：$y=f(x)$。屬性裡可以改成對數或雙對數刻度。
\item 2D 極座標：$r=f(\theta)$，或 $(r,\theta)$ 兩個分量。
\item 2D 隱函數：等式，例如 $x^{2}+y^{2}=4$。
\item 2D 參數式：$(x(t),y(t))$。
\item 2D 保角：複數函數，例如 $\frac{z-1}{z+1}$，畫出方格被函數帶到哪裡。
\item 2D 梯度：$f(x,y)$ 的梯度箭頭。2D 向量場：$[P,Q]$。
\item 2D 不等式：例如 $x+y^{2}<5$，塗滿成立的區域。
\item 2D 常微分方程：把方程和起始值打成一行矩陣，例如 $y'=\sin (xy)$ 和 $y(0)=3$，用數值方法畫解。可以是方程組。
\item 近似積分圖：計算 \TEXTsymbol{>} 微積分 \TEXTsymbol{>} 畫近似積分圖，選左端點、中點、右端點、上和、下和或梯形。
\item 3D 直角座標：$z=f(x,y)$，或參數曲面 $(x(s,t),y(s,t),z(s,t))$。
\item 3D 柱座標：$r=f(\theta ,z)$；球座標：$\rho =f(\theta ,\varphi )$。也可以寫三個分量。
\item 3D 隱函數：等式，例如 $x^{2}+y^{2}+z^{2}=9$。
\item 3D 管狀：沿著曲線 $(x(t),y(t),z(t))$ 的管子，屬性裡可以改半徑。
\item 3D 梯度、3D 向量場 $[P,Q,R]$。
\end{itemize}

\includegraphics{manual/zh-TW/plots.png}

\emph{繪圖 2D > 直角座標 和 繪圖 3D > 直角座標，各用在上面那條公式。}

\includegraphics{manual/zh-TW/plot-props.png}

\emph{繪圖的屬性（選取繪圖按 Ctrl+F5）：公式、顏色、再加一個函數、範圍。}

\section{排版}

\begin{itemize}
\item 排版 PDF（桌面版）：電腦有裝 XeLaTeX 就用它，沒有就用內附的 Tectonic。常用的 TeX 檔已經附在程式裡，沒有網路也能排；用到其他套件時，第一次要上網下載。中文可以排。
\item 導言：排版用的 LaTeX 前言（\TEXTsymbol{\backslash}usepackage 等）。
\item 選項：文件類別、字級、紙張、套件。
\end{itemize}

\section{工具}

\includegraphics{manual/zh-TW/menu-tools.png}

\emph{工具選單。}

\subsection{考卷產生器}

把目前的文件當範本，產生好幾份數字不一樣的考卷和答案。範本是 5.5 的考卷格式（.qiz），用節標題當關鍵字：

\begin{itemize}
\item Exam（或 Quiz、Test）：表示這是考卷範本。Comment：只是註記，不會出現在考卷。
\item Text：原樣放進考卷的文字。Part：分成幾個部分，各有自己的設定。
\item Setup：設定。例如 $a:=\func{rand}(1,9)$（只算一次）、$b=1/a$（每次用到重算）、「Condition: $(a\neq b)$」（不成立就重抽）、「Questions: number, permute」（題目編號、打亂）、「Choices: permute」（選項打亂）、「Seed: 123」（固定亂數）。
\item 亂數：$\func{rand}()$、$\func{rand}(n)$ 是 0 到 $n-1$、$\func{rand}(n,m)$ 是 $n$ 到 $m$、$\func{rand}(\{2,3,5\})$ 從集合挑一個、$\func{randmat}(3,3,1,(1,9))$ 隨機矩陣（第三個數：1 任意、2 對稱、3 反對稱、4 行列式為 1）。
\item Question：一題。Statement：題目。Substatement：小題。Response：選項清單或填空（寫 InputField）。正確的選項後面放 \TEXTsymbol{\backslash}correctchoice（TeX 欄位）。標籤是 fixed 的選項固定放最後。
\item Answer、Solution：只放在答案裡。Variant：同一題的另一種問法，產生時隨機挑一種。
\item 題目裡要代數字的地方用 插入 \TEXTsymbol{>} 考卷公式，例如 $a+b$，產生考卷時才算。
\end{itemize}

產生的考卷開成新的文件，範本還在 視窗 選單。沒做：線上作答和評分、成績資料庫。

\subsection{其他工具}

\begin{itemize}
\item 拼字檢查：英文，用瀏覽器內建的。
\item 文件管理：列出文件用到的圖片和子文件、有沒有找到；打包成 .zip（.sci、5.x .tex、圖片一起）給別人；開啟 .zip 或 5.5 的 .rap 包裝檔。桌面版跟 5.5 一樣可以複製（存檔的文件連同它的檔案複製到別的資料夾）、改名、刪除、清除輔助檔（.aux、.log、.bak……）；刪除的檔案放進資源回收筒，旁邊的每個檔案都可以選擇不要。
\item 自動替換：在公式裡打一串字母（例如 sum10）就換成一個公式。可以整個關掉。
\item 使用者設定：矩陣和表格預設大小、開程式時開哪份文件、存檔留 .bak、公式顏色、介面語言。
\item 匯出設定：HTML 的公式用 MathML 或 MathJax；複製成圖片的解析度。
\item 計算引擎：選 SymPy 或 Mathematica（桌面版、電腦要有裝），並看版本。Lean 設定：Lean 專案資料夾。
\item 計算引擎裡也有會改變答案的設定，和 5.5 的計算設定一樣：解只要主值；對數函數的底從 $e$ 改成 $10$（這時 $\log x$ 是 $\log_{10}x$，計算結果也這樣寫；$\ln$ 還是自然對數）；向量微積分和 $\nabla$ 的基底變數（預設 $x,y,z$）。
\end{itemize}

\includegraphics{manual/zh-TW/user-setup.png}

\emph{工具 > 使用者設定。}

\includegraphics{manual/zh-TW/engine.png}

\emph{工具 > 計算引擎：選引擎，和會改變答案的設定。}

\section{視窗}

可以同時開好幾份文件。視窗選單切換，Ctrl+F6 或 Ctrl+Tab 換下一份。

\section{說明}

\begin{itemize}
\item 快捷鍵（F1）：5.5 的快捷鍵表，再按一次 F1 關掉。
\item 完整說明文件：就是這份文件。
\item 使用說明：一頁的簡單說明。
\item 跟 AI 說話：可以怎麼跟 AI 程式說的例子（見「讓 AI 幫忙」）。
\end{itemize}

\includegraphics{manual/zh-TW/shortcuts.png}

\emph{說明 > 快捷鍵（F1）。}

\includegraphics{manual/zh-TW/guide.png}

\emph{說明 > 使用說明。}

\section{5.5 的檔案}

\begin{itemize}
\item 公式、定理、清單、表格、圖片、繪圖、註腳、引用、索引、子文件、考卷公式都會讀進來。
\item 不認得的 5.x 指令會原樣留成灰底的 TeX 欄位，存檔和匯出都不會掉。
\item 圖片要跟 .tex 放在一起（或 graphics 資料夾）。網頁版讀不到電腦上的圖片，用桌面版，或用 文件管理 開打包檔。
\item 存回 5.x 的 .tex 時，5.5 打得開。
\end{itemize}

\section{讓 AI 幫忙}

桌面版可以讓 Claude、ChatGPT 這類 AI 程式讀 opensw 開著的文件、改文件、計算、排版（透過 MCP）。

\begin{enumerate}
\item 工具 $>$ 連接 AI：勾 Claude 或 ChatGPT / Codex，按「確定」。opensw 會把自己加進那個程式的設定（舊設定先留一份）。把 AI 程式完全關掉再打開。
\item 用自己的話跟它說，例如「讀 opensw 開著的文件，告訴我在講什麼」「改錯字，公式不要動」「算第三條公式裡的積分，答案寫在後面」。說明 $>$ 跟 AI 說話 有更多例子。
\item AI 在改的時候狀態列會顯示，改過的段落有底色。編輯 $>$ 復原 AI 的修改（或狀態列上的按鈕）一次把文件改回原樣。
\end{enumerate}

AI 只能在目前文件的資料夾、文件\textbackslash{}opensw（新文件可以直接存在這裡，不用問）、和 工具 $>$ 使用者設定 允許的資料夾開檔存檔；其他地方 opensw 會先問你。它只能存 .tex 和 .sci。手動設定的方法在原始碼的 docs/AI.md。

裝不了 AI 程式的電腦，或要寫腳本時，可以在命令列用 opensw-cli.exe 下一樣的指令，它在 opensw.exe 旁邊。例如 opensw-cli help 列出指令，opensw-cli save 存開著的文件。指令會傳給開著的編輯器（沒開會自動打開），能開、能存的資料夾也一樣。

\section{關於}

opensw 用 Mozilla Public License 2.0 授權（0.8.0 以前的版本是 MIT）。程式是重新寫的，沒有用 Scientific Word 的程式碼、圖示或說明檔。Scientific Word 是 MacKichan Software 的商標，opensw 跟他們沒有關係。
